% =============================================================================
% Section III: Decoherence from Soft Radiation
% paper/decoherence_framework.tex
% =============================================================================

\section{Decoherence from Soft Radiation}
\label{sec:decoherence}

\subsection{Physical setup}
\label{sec:setup}

We consider a body of charge $q$ (coupling to a scalar field) whose center-of-mass
degree of freedom is placed in a quantum superposition of two spatially separated
worldlines $\gamma_1$ and $\gamma_2$ in the FLRW spacetime $(M, g_{ab})$ described
in \cref{sec:flrw}.  The two worldlines are comoving geodesics that share the same
spatial trajectory up to a proper displacement $d(\tau)\, s^a$, where $s^a$ is a
unit spacelike vector and $d(\tau) \geq 0$ is a real envelope that is switched on
and off \emph{adiabatically}: $d(\tau) \to 0$ for $\tau \to \tau_{\rm on}$ and
$\tau \to \tau_{\rm off}$, with $d(\tau) = d_0 > 0$ during the active interval
$(\tau_{\rm on}, \tau_{\rm off})$ of total proper duration $T =
\tau_{\rm off} - \tau_{\rm on}$.  The adiabatic switching suppresses radiation
associated with the turn-on and turn-off transients~\cite{DSW2022,DSW2025a},
ensuring that the decoherence accumulated during $T$ is unambiguously identified
with the steady-state emission.

The body sources a conformally coupled massless scalar field $\varphi$, satisfying
\be
  \lb( \Box_g - \tfrac{1}{6} R \rb)\, \varphi = S \,,
  \label{eq:wave-eq}
\ee
where $S$ is the scalar charge current.  The key property of \cref{eq:wave-eq} is its
\emph{conformal invariance}: under $g_{ab} \to a^2(\eta)\,\tilde{g}_{ab}$,
\cref{eq:wave-eq} reduces to the massless wave equation $\Box_{\tilde{g}}\chi = 0$
for the rescaled field $\chi = \Scfac(\eta)\,\varphi$~\cite{Wald84}.  This mapping
lies at the heart of the framework developed in
\cref{sec:confvac,sec:deco-formula} below.  The electromagnetic case, which is
likewise conformally invariant in four spacetime dimensions, is treated in
\cref{sec:em}.

The initial state of the field is taken to be the \emph{conformal vacuum}
$|\Psi_0\rangle$ — the state obtained by pulling back the Minkowski vacuum via
the conformal equivalence between the FLRW spacetime and flat Minkowski space
(see \cref{sec:confvac}).  The total initial state of the system is therefore
\be
  |\Phi_0\rangle
  =
  \frac{1}{\sqrt{2}}\lb(\kt{\gamma_1} + \kt{\gamma_2}\rb) \otimes \kt{\Psi_0} \,,
  \label{eq:initial-state}
\ee
where $\kt{\gamma_A}$ denotes the center-of-mass state of the body following
worldline $\gamma_A$.  Soft radiation emitted by the moving charge propagates
outward, and — when a future causal horizon $\mc{H}^+(\gamma)$ is present —
a portion of that radiation escapes beyond the horizon and becomes permanently
inaccessible to any observer along $\gamma$.  The entanglement between the
accessible and inaccessible field degrees of freedom drives the decoherence
of the spatial superposition.


\subsection{The DSW decoherence formula}
\label{sec:deco-formula}

We apply the local decoherence framework of Danielson, Satishchandran, and
Wald~\cite{DSW2022,DSW2023,DSW2025a}.  After tracing over the unobserved field
modes (those beyond the causal horizon), the reduced density matrix of the
body satisfies~\cite{DSW2025a}
\be
  \langle \gamma_1 | \rho_{\rm red}(\tau_{\rm off}) | \gamma_2 \rangle
  =
  e^{-\Gamma/2}\,
  \langle \gamma_1 | \rho_{\rm red}(\tau_{\rm on}) | \gamma_2 \rangle \,,
  \label{eq:rho-off-diag}
\ee
where $\Gamma \geq 0$ is the \emph{decoherence exponent}.  The off-diagonal element
$\langle \gamma_1 | \rho_{\rm red} | \gamma_2 \rangle$ is the quantum coherence
between the two branches; $\Gamma = 0$ means no decoherence, while
$\Gamma \gg 1$ means complete decoherence.  It is convenient to introduce the
\emph{decoherence number}
\be
  \Nnum \equiv \frac{\Gamma}{2}
  = -\ln\abs{\langle \gamma_1 | \rho_{\rm red} | \gamma_2 \rangle} \,,
  \label{eq:Ndef}
\ee
so that $\Nnum \geq 0$ with $\Nnum \gg 1$ indicating full decoherence.

The exact expression for $\Gamma$ in terms of the in-vacuum field
operator $\varphi^{\inc}$ is~\cite{DSW2025a}
\be
  \Gamma
  =
  \langle \Psi_0 |\,
    \lb[\varphi^{\inc}(j_1 - j_2)\rb]^2
  \,|\Psi_0\rangle \,,
  \label{eq:Gamma-exact}
\ee
where $j_A$ is the \emph{classical} scalar current of the body following
path $\gamma_A$, and the smeared field operator is defined by
$\varphi^{\inc}(j) \equiv \int \! j(x)\, \varphi^{\inc}(x)\, \sqrt{-g}\, d^4x$.
Equation~\eqref{eq:Gamma-exact} is exact at leading order in the weak-coupling
expansion (i.e., to order $q^2$) and requires no assumption about the geometry
other than the existence of a well-defined in-vacuum state~\cite{DSW2025a}.
Note that $\Gamma$ depends only on the \emph{difference} of the two source currents;
the individual current of either branch does not enter.  This reflects the physical
picture: a uniformly moving charge does not radiate, and only the relative motion
between the two branches generates the entangling radiation.

\subsubsection*{Dipole approximation}

When the superposition separation $d(\tau)$ is small compared to all relevant
field wavelengths — equivalently, when $d$ is much smaller than the comoving
horizon scale (see \cref{app:dipole}) — the difference of
source currents reduces to a dipole:
\be
  j_1(x) - j_2(x)
  \approx
  q\, d(\tau)\, s^a\, \nabla_a\, \delta^{(3)}\!\lb(\mb{x} - \mb{X}(\tau)\rb) \,,
  \label{eq:dipole}
\ee
where $\mb{X}(\tau)$ denotes the comoving spatial position of Alice's lab (taken
to be the origin, $X^i = 0$), $s^a$ is the unit separation vector, and
$\nabla_a$ acts on the spatial coordinates.  The validity of
\cref{eq:dipole} requires $d \ll \lambda_{\rm min}$ for all modes that contribute
appreciably to $\Gamma$; we verify this condition in \cref{app:dipole}.

Substituting \cref{eq:dipole} into \cref{eq:Gamma-exact} and using the definition
of the Wightman 2-point function
\be
  \Wight(x, x')
  \equiv
  \langle \Psi_0 |\, \varphi^{\inc}(x)\, \varphi^{\inc}(x') \,|\Psi_0\rangle \,,
  \label{eq:Wight-def}
\ee
the decoherence exponent becomes
\begin{align}
  \Gamma
  &=
  q^2
  \int d\tau\!\int d\tau'\;
  d(\tau)\, d(\tau')\;
  s^a s^b\,
  \partial_a \partial_{b'}\,
  \Wight\!\lb(\tau, \mb{X};\; \tau', \mb{X}\rb)
  \bigg|_{\mb{x} = \mb{x}' = \mb{X}} ,
  \label{eq:Gamma-dipole}
\end{align}
where $\partial_a$ and $\partial_{b'}$ denote derivatives with respect to the
spatial coordinates of the first and second arguments of $\Wight$, evaluated at
the worldline $\mb{X}$.  This is the master formula whose evaluation in specific
cosmological backgrounds constitutes the main technical content of
\cref{sec:powerlaw,sec:desitter}.


\subsection{Conformal vacuum and Wightman function}
\label{sec:confvac}

The decoherence formula~\eqref{eq:Gamma-dipole} requires an explicit expression
for the Wightman 2-point function $\Wight(x,x')$ in the physical FLRW spacetime.
We now exploit the conformal invariance of \cref{eq:wave-eq} to obtain this
expression from the known flat-space Wightman function.

\subsubsection*{Conformal vacuum}

In the conformal coordinates $(t, \mb{x})$ introduced in \cref{sec:flrw}, the FLRW
metric takes the form $ds^2 = \Scfac^2(t)\,\eta_{\mu\nu}\, dx^\mu dx^\nu$.  The
conformally rescaled field $\chi \equiv \Scfac(t)\,\varphi$ satisfies the massless
Klein-Gordon equation in the flat metric $\eta_{\mu\nu}$, so its mode decomposition
is identical to that of a massless field in Minkowski space.  We define the
\emph{conformal vacuum} $|\Psi_0\rangle$ to be the state in which these Minkowski
modes are in their ground state — that is, $|\Psi_0\rangle$ is the pullback of the
Minkowski vacuum under the conformal map.

For de Sitter spacetime, the conformal vacuum coincides with the well-known
Bunch-Davies state~\cite{BunchDavies1978,Allen1985} (see \cref{app:bunch-davies}
for a direct comparison of the Wightman functions).  More generally, the conformal
vacuum is the natural initial state for any spatially flat FLRW spacetime with
conformally invariant matter: it is the unique Hadamard state that is invariant
under the Euclidean symmetries of the spatial slices and that reduces to the
Minkowski vacuum in the flat-space limit $\Scfac \to \mathrm{const}$.

\subsubsection*{Wightman function in the conformal vacuum}

The 2-point function of the original (unrescaled) field $\varphi$ in the conformal
vacuum follows immediately from the flat-space Wightman function $W_{\rm Mink}$
via the conformal weight:
\be
  \Wight(x_1, x_2)
  =
  \frac{1}{\Scfac(t_1)\,\Scfac(t_2)}\,
  W_{\rm Mink}(\tilde{x}_1, \tilde{x}_2) \,,
  \label{eq:Wight-conformal-relation}
\ee
where $\tilde{x}$ denotes the same coordinate point viewed in the conformally
related Minkowski spacetime.  Since the Minkowski massless scalar Wightman function
is
$W_{\rm Mink}(\tilde{x}_1,\tilde{x}_2)
  = -[4\pi^2 \tilde{\sigma}_\varepsilon(\tilde{x}_1,\tilde{x}_2)]^{-1}$
with $\tilde{\sigma}_\varepsilon = -(t_1-t_2-i\varepsilon)^2+|\mb{x}_1-\mb{x}_2|^2$
(in our signature conventions), \cref{eq:Wight-conformal-relation} gives
\be
  \Wight(x_1, x_2)
  =
  \frac{-1}{4\pi^2\, \Scfac(t_1)\, \Scfac(t_2)\,
    \tilde{\sigma}_\varepsilon(x_1, x_2)} \,,
  \label{eq:Wight-FLRW}
\ee
where the regularized squared geodesic distance in the conformal Minkowski frame is
\be
  \tilde{\sigma}_\varepsilon(x_1, x_2)
  =
  -\lb(t_1 - t_2 - i\varepsilon\rb)^2
  + \abs{\mb{x}_1 - \mb{x}_2}^2 \,,
  \quad \varepsilon \to 0^+ \,.
  \label{eq:sigma-tilde}
\ee
Here $t = \int d\tau/\Scfac(\tau)$ is the conformal time coordinate established
in \cref{sec:flrw}, and the $i\varepsilon$ prescription implements the standard
Wightman boundary condition (positive-frequency modes dominate).
Equations~\eqref{eq:Wight-FLRW}--\eqref{eq:sigma-tilde} agree with the
conformal-vacuum Wightman function quoted in
Refs.~\cite{DSW2025a,BunchDavies1978,Allen1985}.

\subsubsection*{Role of the causal horizon}

The decoherence exponent~\eqref{eq:Gamma-dipole} involves a double integral over
the conformal-time history of the superposition.  The temporal support of the
integrand is $t \in [t_{\rm on}, t_{\rm off}]$, where
$t_{\rm on} = \int_{\tau_*}^{\tau_{\rm on}} d\tau/\Scfac(\tau)$ and
$t_{\rm off} = \int_{\tau_*}^{\tau_{\rm off}} d\tau/\Scfac(\tau)$.

The presence or absence of a future causal horizon is determined by whether the
conformal time range is finite or infinite as $\tau_{\rm off} \to \infty$:
\begin{itemize}[leftmargin=2em]
  \item \emph{Horizon present} ($p > 1$ or de Sitter): the integral
    $t_{\infty} \equiv \int_{\tau_*}^{\infty} d\tau/\Scfac(\tau)$ converges.
    The conformal-time range $[t_{\rm on}, t_{\rm off}]$ is bounded above by
    $t_\infty < \infty$.  The support of $d(\tau)$ is confined to this finite
    interval, which acts as a \emph{physical infrared (IR) cutoff} on the
    field modes contributing to $\Gamma$.

  \item \emph{No horizon} ($p \leq 1$, $p \neq 0$): the conformal-time range is
    unbounded, $t_\infty = \infty$.  The IR modes then contribute a piece that
    saturates to a constant independent of $T$, so $\Nnum \to \mathrm{const}$
    as $T \to \infty$ — consistent with the absence of any ongoing loss of
    information through a horizon.
\end{itemize}
This dichotomy is the central mechanism by which the causal horizon converts a
potentially IR-divergent flat-space formula into a finite, physically meaningful
decoherence rate.  The explicit computation confirming this picture for the
power-law family and for de Sitter is carried out in
\cref{sec:powerlaw,sec:desitter}, respectively.

%% CITATIONS NEEDED (for gpd-bibliographer):
%% - DSW2022        -- Danielson, Satishchandran, Wald (2022a) IJMPD 31, 2241003
%% - DSW2023        -- Danielson, Satishchandran, Wald (2023) PRD 108, 025007
%% - DSW2025a       -- Danielson, Satishchandran, Wald (2025a) PRD 111, 025014
%% - BunchDavies1978 -- Bunch & Davies (1978) Proc. R. Soc. Lond. A 360, 117
%% - Allen1985      -- Allen (1985) PRD 32, 3136
%% - Wald84         -- Wald, General Relativity (1984) [for conformally coupled scalar R/6 term]
